\begin{ejer}
Demuestra que para cualquier ideal $J\subseteq S^{-1}A$ se tiene
\(
J=(\varphi^{-1}(J))^e.
\)
Deduce que la biyección $\operatorname{Spec}A_f\cong D(f)$ es un homeomorfismo.
\end{ejer}
\begin{sol}
En primer lugar:
\[
\frac{a}{s}\in J
\quad%
\underset{%
    \mathclap{\substack{%
    \big\uparrow\\\text{ya que }\frac{s}{1}\text{ es invertible}%
    }}
}{\iff}%
\quad%
\varphi(a)=\frac{a}{1}=\frac{s}{1}\frac{a}{s}\in J
\quad%
\underset{%
    \mathclap{\substack{%
    \big\uparrow\\\text{definición}%
    }}
}{\iff}%
\quad%
a\in\varphi^{-1}(J)
\quad%
\underset{%
    \mathclap{\substack{%
    \big\uparrow\\\text{definición}%
    }}
}{\iff}%
\quad%
\frac{a}{s}\in(\varphi^{-1}(J))^e
\]
Esto prueba $J=(\varphi^{-1}(J))^e$. Para $A_f$ el conjunto multiplicativo
es $S=\set*{f^n\given n\en\mathbb{N}_0}$. El teorema de correspondencia para
localizaciones da una biyección
\[
\Phi\colon\operatorname{Spec}A_f%
\;\xlongrightarrow{\cong}\;%
\set*{\mathfrak{p}\en\operatorname{Spec}A\given \mathfrak{p}\cap S=\varnothing}%
\;=\;%
D(f)
\]
ya que $\mathfrak{p}\cap S=\varnothing$ si y solo si $f\notin\mathfrak{p}$.

Por definición, $\Phi\colon\operatorname{Spec}A_f\to\operatorname{Spec}A$ es
inducida por $\varphi\colon A\to A_f$, así que el Ejercicio 1.6 implica que
es continua. Más aún, sea $J\subseteq A_f$ un ideal y pongamos
$I=\varphi^{-1}(J)$. Por lo probado al inicio, $J=I^e$. Si
$\mathfrak{p}\en D(f)$, el primo de $A_f$ que le corresponde mediante
$\Phi^{-1}$ es $\mathfrak{p}^e$. Por tanto,
\begin{align*}
\mathfrak{p}\en\Phi\bigl(V_{A_f}(J)\bigr)
&\quad\Longleftrightarrow\quad
\mathfrak{p}\en D(f)\ \wedge\ \mathfrak{p}^e\en V_{A_f}(J)\\
&\quad\Longleftrightarrow\quad
\mathfrak{p}\en D(f)\ \wedge\ \mathfrak{p}^e\supseteq J=I^e\\
&\quad\Longleftrightarrow\quad
\mathfrak{p}\en D(f)\ \wedge\ \mathfrak{p}\supseteq I\\
&\quad\Longleftrightarrow\quad
\mathfrak{p}\en D(f)\cap V_A(I)
\end{align*}
Esto prueba que
$\Phi\bigl(V_{A_f}(J)\bigr)=D(f)\cap V_A\bigl(\varphi^{-1}(J)\bigr)$,
por lo que $\Phi$ manda cerrados de $\operatorname{Spec}A_f$ en cerrados de
$D(f)$. Así, $\Phi$ es cerrada. Como $\Phi$ es continua, biyectiva y cerrada,
es homeomorfismo.
\end{sol}